\documentclass[11pt]{article}

\usepackage[utf8]{inputenc}
\usepackage[spanish]{babel}
\usepackage[margin=2.5cm]{geometry}
\usepackage{listings}
\usepackage{xcolor}
\usepackage{amsmath}
\usepackage{graphicx}
\usepackage{float}

% Estilo
\lstdefinestyle{pythonstyle}{
    language=Python,
    backgroundcolor=\color{gray!10},
    basicstyle=\ttfamily\footnotesize,
    keywordstyle=\color{blue}\bfseries,
    commentstyle=\color{gray!70!black}\itshape,
    stringstyle=\color{orange},
    numbers=left,
    numberstyle=\tiny\color{gray},
    stepnumber=1,
    numbersep=8pt,
    breaklines=true,
    frame=single,
    showstringspaces=false,
    tabsize=4
}
\lstset{style=pythonstyle}

\title{Seguridad Informática - Práctica 2 -- Algoritmo RSA}
\author{Mauricio A Sonda Cahuich - 66611}
\date{\today}

\begin{document}

\maketitle

\section{¿Qué hice?}

Para esta práctica programé en Python el algoritmo RSA para descifrar un
mensaje que nos dieron, usando los primos $p = 67$ y $q = 83$. La idea era
calcular $n$, $\varphi(n)$, encontrar el cuarto valor de $e$ que sirviera
como clave pública, sacar su $d$ correspondiente (clave privada) y con eso
descifrar el mensaje.

Lo que hice diferente fue no solo contar los candidatos de $e$ hasta llegar
al cuarto, sino que hice que el programa \textbf{probara cada uno} intentando
descifrar el mensaje real, para comprobar cuál de verdad funciona y así tener
una justificación más sólida (y no solo "porque sí es el cuarto").

\section{Datos de inicio}

\begin{align*}
p &= 67 \\
q &= 83
\end{align*}

Mensaje cifrado:
\begin{center}
1058, 2597, 4955, 4201, 3162, 4343, 2754, 2497, 336, 2597
\end{center}

\section{Paso 1: Sacar $n$}

Esto es lo más fácil, nada más se multiplican los primos:

\[
n = p \times q = 67 \times 83 = 5561
\]

\section{Paso 2: Sacar $\varphi(n)$}

\[
\varphi(n) = (p-1)(q-1) = 66 \times 82 = 5412
\]

En el código nada más los calculo y los imprimo para tener registro:

\begin{lstlisting}
n = p * q
phi = (p - 1) * (q - 1)

print("n = " + str(n))
print("phi(n) = " + str(phi) + "\n")
\end{lstlisting}

\section{Paso 3: Buscar el $e$ correcto}

Aquí es donde le puse más trabajo. Sabemos que $e$ tiene que
cumplir:

\[
1 < e < \varphi(n) \quad \text{y} \quad \mathrm{mcd}(e, \varphi(n)) = 1
\]

Entonces hice un while que va probando números desde el 2 hacia arriba, y
cada vez que encuentra uno que cumple lo del mcd, lo cuenta como candidato.
Pero en vez de solo quedarme con el cuarto y ya, hice que el programa
\textbf{intentara descifrar el mensaje con cada candidato}, y si el
resultado eran puros caracteres raros/no imprimibles, lo descartaba y seguía
con el siguiente.

\begin{lstlisting}
e = 2
while e < phi and not encontrado:

    if mcd(e, phi) != 1:
        e += 1
        continue  # no cumple, ni lo cuento

    contador += 1
    d = calcularD(e, phi)

    print("Candidato #" + str(contador) + ": e = " + str(e) + ", d = " + str(d))

    intento = descifrarMensaje(mensajeCifrado, d, n)

    if esTextoValido(intento):
        print("Este si funciona! usamos e = " + str(e))
        encontrado = True
    else:
        print("Este no sirve, el texto sale con basura.")

    e += 1
\end{lstlisting}

Para saber si el texto "sirve" hice una función que checa que todos los
caracteres estén en el rango normal de ASCII imprimible (del 32 al 126,
o sea letras, números y símbolos normales):

\begin{lstlisting}
def esTextoValido(texto):
    for c in texto:
        if ord(c) < 32 or ord(c) > 126:
            return False  # salio un caracter raro, este e no sirve
    return True
\end{lstlisting}

\subsection{¿Por qué el cuarto candidato es el correcto?}

Al correr el programa, los primeros candidatos que cumplen lo del mcd son
$5$, $7$, $13$ y $17$. Cuando los probé uno por uno descifrando el mensaje:

\begin{itemize}
    \item Con $e = 5$ el mensaje sale puro símbolo raro, no sirve.
    \item Con $e = 7$ pasa lo mismo, sale basura.
    \item Con $e = 13$ otra vez sale texto sin sentido.
    \item Con $e = 17$ (el cuarto) el mensaje \textbf{sí sale legible}, con
    letras normales.
\end{itemize}

Entonces ahí queda comprobado: $e = 17$ no solo es el cuarto número que
cumple la condición matemática, además es el que de verdad
descifra bien el mensaje. Eso es la prueba de que es el correcto.

\section{Paso 4: Sacar $d$ (clave privada)}

$d$ tiene que cumplir que:

\[
e \times d \equiv 1 \pmod{\varphi(n)}
\]

En vez de usar el algoritmo extendido de Euclides a mano, hice que el
programa lo fuera probando con un while hasta encontrar el valor que
cumpliera:

\begin{lstlisting}
def calcularD(e, phi):
    d = 1
    while (d * e) % phi != 1:
        d += 1  # le sigo subiendo hasta que de resto 1
    return d
\end{lstlisting}

\section{Paso 5: Descifrar el mensaje}

Para descifrar cada número del mensaje se usa la fórmula:

\[
m = c^{d} \bmod n
\]

Lo hice con un for que va multiplicando y sacando el módulo poco a poco
(en vez de elevar todo de un jalón, porque los números se hacen enormes):

\begin{lstlisting}
def descifrar(c, d, n):
    m = 1
    # aqui se hace (c elevado a la d) mod n, pero poco a poco
    for i in range(d):
        m = (m * c) % n
    return chr(m)
\end{lstlisting}

\section{Resultados}

\begin{itemize}
    \item \textbf{Clave pública:} $(e, n) = (17,\ 5561)$
    \item \textbf{Clave privada:} $(d, n)$, la calcula el programa
    \item \textbf{Mensaje cifrado:} 1058, 2597, 4955, 4201, 3162, 4343, 2754,
    2497, 336, 2597
    \item \textbf{Mensaje descifrado:} HOLA MUNDO
\end{itemize}

\section{Evidencias}

\begin{itemize}
    \item \textbf{Evidencia:} Captura con el resultado final en terminal
\end{itemize}

% Imagen
\begin{figure}[H]
    \centering
    \includegraphics[width=0.9\textwidth]{terminalRSA.png}
    \caption{Consola mostrando el calculo e intentos.}
\end{figure}

\section{Conclusión}

Con este programa no nada más se sacan las claves y se descifra el mensaje,
sino que traté de hacer visible de manera clara cómo se llega al resultado
final, y el porque el cuarto candidato de $e$ es el
correcto: se prueban los candidatos válidos hasta que uno
logra descifrar el mensaje de forma legible. Eso deja comprobado, con
evidencia del propio programa, que la elección de $e = 17$ es la correcta.

\end{document}